\documentclass[a4paper, 11pt]{article}
\usepackage{xeCJK}
\usepackage{geometry}
\usepackage{titlesec}

\geometry{margin=1in}
\setCJKmainfont{MS Mincho}

% \title{進捗報告}
% \author{コレドール・パブロ}
% \date{２０２５年１０月１０日}

\titlespacing*{\section}{0pt}{2ex}{1ex}
\titlespacing*{\subsection}{0pt}{1ex}{0.8ex}
\titlespacing*{\subsubsection}{0pt}{1ex}{0.5ex}

\begin{document}

\begin{center}
    {\LARGE \bfseries 進捗報告} \\ % Title: Large and bold
    \vspace{0.3em} % A small vertical space
    Juan Pablo Corredor \\ % Author
    ２０２５年１１月１４日 % Date
\end{center}
\vspace{0.5em} % Adjust space between title block and main text

\section*{近況}
% Write your recent findings here.
\begin{itemize}
    \item アニメ科コラボのBGMを完成した。ミックスセッションに参加した。
    \item 接種、プレゼン等は無事にした。
\end{itemize} 

\section*{やったこと}
\begin{itemize} 
    \item 分析方法をテスト中、そのため、GUI、アプリ、Batch分析の機能を開発した。 
    \item サイン分析の結果が大分改良。 
    \item ノイズ分析はなかなか進まない。FFT、周波数帯域、エネルギーの数学の難易度が高く、使用している式を確認した上でも、エネルギーの値が正しく分析できない。\cite{sitrano}での分析方法はSTFT-ISTFTすなわち、スペクトルグラムから音源に戻すアルゴリズムを使用。
\end{itemize}

\section*{考察}

出力していたCSVファイルのサイズは大きくて、データ扱いは非効率的な作用になるため、分析結果はバイナリーファイルに書き込んでいる（独自な拡張子で）。

FFTウィンドウ（一周期、周期ループ、直接音源）の比較・評価方法を探している。

また、機械学習モデルの出力を線形補間と比較する予定であるため、バイナリー倍音データを補間し、音源を生成するプログラムも開発中。

聴覚実験が必要になる可能性もあると考えている。

現在研究計画：

倍音データの線形補間。

機械学習モデルの開発を進んで、完了したら、補間。

補間に基づき生成される音源の比較・評価。

プラグインを開発し、最も精度が高いアルゴリズムを総合。


\section*{やること}
% Outline future steps here.
\begin{itemize}
    \item 開発を続ける。
    \item ノイズの分析を解決するように頑張る。
\end{itemize}

% Add your bibliography here.
\bibliography{../../ref}
\bibliographystyle{apalike}

\end{document}